% =============================================================================
% 08-cpu-scheduling-basics-notes.tex — Lecture Notes for Lecture 8:
% CPU Scheduling — Concepts, Criteria, FCFS and SJF
%
% DERIVED ARTIFACT. Source of truth: Slides/CS502/08-cpu-scheduling-basics.tex.
% Generated by /lecture-notes; do not hand-edit content independently of the
% Beamer source without re-running /qa-notes afterward (single-source-of-truth.md).
% Structured as a textbook chapter (numbered sections/figures/examples,
% topic-organized rather than a frame-by-frame transcript).
%
% Every Gantt chart and CT/TAT/WT/RT table below is the same generated code
% as the deck (verified simulator; FCFS, non-preemptive SJF, SRTF; no ties in
% any instance). Exercise instances were verified with the same simulator.
%
% Anchor reading: Silberschatz, Galvin & Gagne, Operating System Concepts,
%                 8th ed. — Ch. 5 §5.1 (p. 183), §5.2 (p. 187), §5.3.1 (p. 188),
%                 §5.3.2 (p. 189), page-verified via
%                 master_supporting_docs/CS502/supporting_books/SilberschatzOSC/index.md;
%                 Stallings, Operating Systems: Internals and Design Principles,
%                 6th ed. — Ch. 9 §9.1 (p. 406), §9.2 (p. 410), page-verified via
%                 master_supporting_docs/CS502/supporting_books/StallingsOS/index.md.
%
% Compile with: cd Notes/CS502 &&
%   TEXINPUTS=../../Preambles:$TEXINPUTS xelatex -interaction=nonstopmode 08-cpu-scheduling-basics-notes.tex  (x3)
% =============================================================================

\documentclass[11pt]{article}
\usepackage[margin=1in]{geometry}
\usepackage{array}
\input{../../Preambles/header}
\coursecode{CS502}

% Chapter-style numbering: this is Lecture 8's chapter.
\renewcommand{\thesection}{8.\arabic{section}}
\renewcommand{\thefigure}{8.\arabic{figure}}
\renewcommand{\thetable}{8.\arabic{table}}

\newtheorem{example}{Example}
\renewcommand{\theexample}{8.\arabic{example}}

\newenvironment{definitionbox}[1]{%
  \par\noindent\textbf{Definition (#1).}\ \itshape
}{\par\vspace{0.3em}}

\setlength{\emergencystretch}{2em}

\title{CPU Scheduling\\\large Lecture Notes --- Lecture 8\\
  \normalsize Concepts, Criteria, FCFS and SJF}
\author{Auqib Hamid Lone\\[0.3em]
  \emph{Department of Computer Science \& Engineering}, \emph{GCET Ganderbal}}
\date{\today}

\begin{document}
\maketitle

\section{Five Jobs, One CPU}

Five students submit compile jobs to the department's lab server within five
milliseconds of one another. The server has a single CPU. Table~\ref{tab:lab}
lists, for each job, when it arrives and how long a CPU burst it needs. The
question that organizes this chapter is the one the instructor faces at the
console: \emph{in what order should the jobs run, and how would we even decide
that one order is better than another?}

\begin{table}[h]
\centering
\begin{tabular}{lccccc}
  \toprule
  Process & \texttt{P1} & \texttt{P2} & \texttt{P3} & \texttt{P4} & \texttt{P5} \\
  \midrule
  Arrival time (ms) & 0 & 1 & 2 & 3 & 5 \\
  CPU burst $\tau$ (ms) & 6 & 8 & 2 & 5 & 4 \\
  \bottomrule
\end{tabular}
\caption{The lab-server workload, the running example of this chapter.}
\label{tab:lab}
\end{table}

We will schedule these five jobs three ways --- first-come first-served (FCFS),
shortest-job-first (SJF), and shortest-remaining-time-first (SRTF) --- and
measure each schedule with the same yardsticks. Four questions frame the
chapter:

\begin{enumerate}
  \item \key{Why schedule at all?} The CPU--I/O burst cycle, the scheduler, the
        dispatcher, and preemption (Section~\ref{sec:concepts}).
  \item \key{What makes a schedule good?} Five criteria and the formulas for
        turnaround, waiting, and response time (Section~\ref{sec:criteria}).
  \item \key{What does the simplest policy do?} FCFS and the convoy effect
        (Section~\ref{sec:fcfs}).
  \item \key{Can we do better?} SJF, its optimality, burst prediction, and the
        preemptive form SRTF (Sections~\ref{sec:sjf}--\ref{sec:srtf}).
\end{enumerate}

\noindent Lecture~6 introduced the short-term scheduler as a box in the
process-state diagram. This chapter opens that box and asks what policy it
should run.

% -----------------------------------------------------------------------------
\section{Basic Concepts}
\label{sec:concepts}

\noindent\emph{Reading:} Silberschatz et al., Ch.~5 \S5.1, pp.~183--187
\cite{SilberschatzOSC_operating_system_concepts}; Stallings, Ch.~9 \S9.1,
p.~406 \cite{StallingsOS_internals_design_principles}.

\subsection*{The CPU--I/O burst cycle}

A process does not use the CPU continuously. It alternates between computing
and waiting, and the waiting is what makes scheduling worthwhile: while one
process waits for a disk or a keyboard, another can use the CPU.

\begin{definitionbox}{CPU burst and I/O burst}
A CPU burst is an interval during which a process executes on the CPU; an I/O
burst is an interval during which it waits for an I/O operation. Execution is a
cycle of the two, beginning and ending with a CPU burst.
\end{definitionbox}

\begin{figure}[h]
\centering
\begin{tikzpicture}[every node/.style={font=\small}]
    % Coordinate map: six boxes on y=0, centres x = -5.5, -3.3, -1.1, 1.1, 3.3, 5.5;
    % each box 1.9cm wide (gap 0.3cm), extent x = -6.45..6.45. CPU boxes use
    % dag-node, I/O boxes dag-node with gold border, white fill. Time arrow on y=-0.9, x = -6.45..6.45.
    \node[dag-node, minimum width=1.9cm, minimum height=0.8cm] (c1) at (-5.5,0) {CPU burst};
    \node[dag-node, draw=primary-gold, fill=white, minimum width=1.9cm, minimum height=0.8cm] (i1) at (-3.3,0) {I/O burst};
    \node[dag-node, minimum width=1.9cm, minimum height=0.8cm] (c2) at (-1.1,0) {CPU burst};
    \node[dag-node, draw=primary-gold, fill=white, minimum width=1.9cm, minimum height=0.8cm] (i2) at (1.1,0) {I/O burst};
    \node[dag-node, minimum width=1.9cm, minimum height=0.8cm] (c3) at (3.3,0) {CPU burst};
    \node[dag-node, draw=primary-gold, fill=white, minimum width=1.9cm, minimum height=0.8cm] (i3) at (5.5,0) {\ldots};
    \draw[observed-edge] (-6.45,-0.9) -- node[below, font=\scriptsize] {time} (6.45,-0.9);
  \end{tikzpicture}
\caption{The CPU--I/O burst cycle: CPU bursts (blue) alternate with I/O bursts
(gold outline) along the time axis --- cf.\ Silberschatz et al., Ch.~5 \S5.1,
p.~183 \cite{SilberschatzOSC_operating_system_concepts}.}
\label{fig:burst}
\end{figure}

Figure~\ref{fig:burst} reads left to right: a CPU burst, then an I/O burst, then
another CPU burst, and so on until a final CPU burst ends in termination.
Measured burst lengths are dominated by many short bursts and a few long ones.
That distribution separates two kinds of program. An \key{I/O-bound} program has
many short CPU bursts; a \key{CPU-bound} program has a few long ones. A good
scheduler must serve both.

\subsection*{The CPU scheduler and the ready queue}

\begin{definitionbox}{CPU scheduler}
The CPU scheduler (short-term scheduler) selects one process from the processes
in memory that are ready to execute and allocates the CPU to it.
\end{definitionbox}

The ready queue holds PCBs, not programs, and it is \emph{not} necessarily a
FIFO queue: it can be a FIFO queue, a priority queue, a tree, or an unordered
list. The \emph{policy} is the rule that picks from it, and FCFS, SJF, and SRTF
are three such rules. Whenever one process waits for I/O, the scheduler gives
the CPU to another --- this is how multiprogramming keeps the CPU busy. The
long-term and medium-term schedulers of Lecture~6 decide \emph{who is in
memory}; the CPU scheduler of this chapter decides \emph{who runs}.

\subsection*{Preemptive and non-preemptive scheduling}

A scheduling decision may be needed at four kinds of moment, namely when a
process:

\begin{enumerate}
  \item switches from running to waiting (an I/O request, or \texttt{wait});
  \item switches from running to ready (a timer interrupt);
  \item switches from waiting to ready (I/O completion);
  \item terminates.
\end{enumerate}

\noindent In cases~1 and~4 the running process gives up the CPU itself, so the
scheduler has no choice but to pick another process. In cases~2 and~3 it has a
choice.

\begin{definitionbox}{Non-preemptive and preemptive scheduling}
Under non-preemptive (cooperative) scheduling, decisions are made only in
cases~1 and~4: once a process has the CPU, it keeps it until it blocks or
exits. Under preemptive scheduling, decisions are also made in cases~2 and~3:
the kernel may take the CPU away from a running process.
\end{definitionbox}

Preemption has a price. If the kernel can stop a process at any moment, it may
stop it in the middle of updating shared data, and the kernel must then protect
that data from a half-finished update --- the problem the concurrency unit of
this course takes up.

\begin{example}[Which transitions force a decision]
A process finishes its CPU burst and issues a disk \texttt{read}. Must a
non-preemptive scheduler make a decision, and what about when the \texttt{read}
later completes?

\emph{Solution.} Issuing the \texttt{read} is case~1 (running $\to$ waiting):
the CPU is now free, so the scheduler must pick some ready process. The
completion of the \texttt{read} is case~3 (waiting $\to$ ready). A
non-preemptive scheduler lets the currently running process keep the CPU; the
newly ready process simply joins the ready queue. Only a preemptive scheduler
may reconsider at case~3.
\end{example}

\subsection*{The dispatcher}

\begin{definitionbox}{Dispatcher and dispatch latency}
The dispatcher is the module that gives control of the CPU to the process the
scheduler selected. The time it takes to stop one process and start another is
the dispatch latency.
\end{definitionbox}

The dispatcher does three things: it switches context (saving one PCB and
loading another, as in Lecture~6), it switches to user mode, and it jumps to the
saved program counter in the new process. It runs on \emph{every} switch, so a
policy that switches more often pays the dispatch latency more often. We will
measure exactly that cost on the lab-server jobs in Section~\ref{sec:compare}.

% -----------------------------------------------------------------------------
\section{Scheduling Criteria}
\label{sec:criteria}

\noindent\emph{Reading:} Silberschatz et al., Ch.~5 \S5.2, pp.~187--188
\cite{SilberschatzOSC_operating_system_concepts}.

Before comparing policies we need yardsticks: ``better'' means nothing until we
say better for whom, measured how. Table~\ref{tab:criteria} lists the five
standard criteria.

\begin{table}[h]
\centering
\renewcommand{\arraystretch}{1.25}
\begin{tabular}{>{\raggedright\arraybackslash}p{0.24\linewidth}>{\raggedright\arraybackslash}p{0.48\linewidth}l}
  \toprule
  \textbf{Criterion} & \textbf{Meaning} & \textbf{Goal} \\
  \midrule
  CPU utilization & fraction of time the CPU does useful work & maximize \\
  Throughput & processes completed per unit time & maximize \\
  Turnaround time (TAT) & submission to completion, per process & minimize \\
  Waiting time (WT) & total time spent in the ready queue & minimize \\
  Response time (RT) & submission to the \emph{first} time on the CPU & minimize \\
  \bottomrule
\end{tabular}
\caption{The five scheduling criteria (Silberschatz et al., Ch.~5 \S5.2).}
\label{tab:criteria}
\end{table}

Two observations shape everything that follows. First, the scheduling policy
directly controls only the \key{waiting time}: it does not change how long a
burst computes, only how long a process sits in the ready queue before and
between its turns. Second, we usually optimize the \emph{average} of a
criterion, but interactive systems also care about the \emph{variance} of
response time: a predictable response is often worth more than a slightly
faster but erratic one.

\subsection*{The three per-process formulas}

\begin{definitionbox}{TAT, WT, RT}
For process \texttt{Pi} with arrival time $\mathrm{AT}_i$, CPU burst $\tau_i$,
first-dispatch time $\mathrm{ST}_i$, and completion time $\mathrm{CT}_i$ (one
CPU burst per process),
\[
  \mathrm{TAT}_i = \mathrm{CT}_i - \mathrm{AT}_i, \qquad
  \mathrm{WT}_i = \mathrm{TAT}_i - \tau_i, \qquad
  \mathrm{RT}_i = \mathrm{ST}_i - \mathrm{AT}_i .
\]
\end{definitionbox}

The waiting-time formula follows from the turnaround: between arrival and
completion a process is either running (for exactly $\tau_i$ in total) or ready
and not running, so the second part is $\mathrm{TAT}_i - \tau_i$. Hence WT
counts \emph{all} time ready but not running, including time spent back in the
queue after a preemption. For a non-preemptive policy the process runs once,
without interruption, starting at $\mathrm{ST}_i$; then
$\mathrm{CT}_i = \mathrm{ST}_i + \tau_i$ and so $\mathrm{WT}_i = \mathrm{RT}_i$.
Finally, idle CPU time is not anybody's waiting time: when the CPU is idle, no
process is ready.

Gantt charts in this chapter draw time left to right; each bar is labelled with
the process that runs, idle intervals are shaded grey, and arrivals are marked
by gold arrows below the axis.

\begin{example}[Metrics from a schedule with an idle gap]
\label{ex:idle}
Jobs \texttt{P1} $(0,3)$, \texttt{P2} $(5,4)$, \texttt{P3} $(6,2)$, given as
(AT, $\tau$) in ms, are scheduled FCFS. Find every metric and the CPU
utilization.

\emph{Solution.} \texttt{P1} arrives at 0 and runs from 0 to 3. At 3 nothing
else has arrived, so the CPU idles from 3 to 5. \texttt{P2} arrives at 5 and
starts at once, running 5 to 9. \texttt{P3} arrived at 6 and has been waiting;
it runs 9 to 11. The chart is Figure~\ref{fig:idle}.

\noindent Applying the formulas: \texttt{P1}: $\mathrm{TAT}_1 = 3-0 = 3$,
$\mathrm{WT}_1 = 3-3 = 0$. \texttt{P2}: $\mathrm{TAT}_2 = 9-5 = 4$,
$\mathrm{WT}_2 = 4-4 = 0$ (it starts at its arrival). \texttt{P3}:
$\mathrm{TAT}_3 = 11-6 = 5$, $\mathrm{WT}_3 = 5-2 = 3$. Because FCFS is
non-preemptive, each RT equals the WT.

\begin{center}
{\footnotesize\begin{tabular}{lrrrrrr}
\toprule
\textbf{Proc.} & \textbf{AT} & $\boldsymbol{\tau}$ & \textbf{CT} & \textbf{TAT} & \textbf{WT} & \textbf{RT} \\
\midrule
\texttt{P1} & 0 & 3 & 3 & 3 & 0 & 0 \\
\texttt{P2} & 5 & 4 & 9 & 4 & 0 & 0 \\
\texttt{P3} & 6 & 2 & 11 & 5 & 3 & 3 \\
\midrule
\multicolumn{4}{@{}l}{\textbf{Average}} & 4 & 1 & 1 \\
\bottomrule
\end{tabular}}
\end{center}

\noindent The CPU is busy for $3+4+2 = 9$ of the 11 ms, so utilization over
$[0,11]$ is $9/11 \approx 81.8\%$.
\end{example}

\begin{figure}[h]
\centering
\begin{tikzpicture}[x=0.75cm,y=1cm]
  % Gantt: I under FCFS (generated from the verified simulator); time axis left to right,
  % bars y=0..0.6, time labels below y=0, arrival arrows y=-0.95..-0.5, arrival names below y=-0.95
  \draw[draw=primary-blue, fill=light-bg] (0,0) rectangle (3,0.6);
  \node[font=\scriptsize] at (1.5,0.3) {\texttt{P1}};
  \draw[draw=neutral, fill=neutral!12] (3,0) rectangle (5,0.6);
  \node[font=\tiny, text=neutral] at (4,0.3) {idle};
  \draw[draw=primary-blue, fill=light-bg] (5,0) rectangle (9,0.6);
  \node[font=\scriptsize] at (7,0.3) {\texttt{P2}};
  \draw[draw=primary-blue, fill=light-bg] (9,0) rectangle (11,0.6);
  \node[font=\scriptsize] at (10,0.3) {\texttt{P3}};
  \node[font=\tiny, below] at (0,0) {0};
  \node[font=\tiny, below] at (3,0) {3};
  \node[font=\tiny, below] at (5,0) {5};
  \node[font=\tiny, below] at (9,0) {9};
  \node[font=\tiny, below] at (11,0) {11};
  \draw[-{Latex[length=1.2mm]}, draw=primary-gold] (0,-0.95) -- (0,-0.5);
  \node[font=\tiny, text=primary-gold, below] at (0,-0.95) {P1};
  \draw[-{Latex[length=1.2mm]}, draw=primary-gold] (5,-0.95) -- (5,-0.5);
  \node[font=\tiny, text=primary-gold, below] at (5,-0.95) {P2};
  \draw[-{Latex[length=1.2mm]}, draw=primary-gold] (6,-0.95) -- (6,-0.5);
  \node[font=\tiny, text=primary-gold, below] at (6,-0.95) {P3};
\end{tikzpicture}
\caption{FCFS schedule of Example~\ref{ex:idle}, with an idle gap from 3 to 5.}
\label{fig:idle}
\end{figure}

\subsection*{Utilization and throughput include the dispatcher}

\begin{example}[Dispatch latency on the lab server]
\label{ex:util}
The lab-server jobs have total burst $6+8+2+5+4 = 25$~ms. FCFS dispatches each
job once, so there are 5 dispatches. With a dispatch latency of 0.5~ms per
dispatch, the CPU spends $25 + 5 \times 0.5 = 27.5$~ms to finish all five jobs,
of which 25~ms is useful work. Therefore
\[
  \text{utilization} = \frac{25}{27.5} \approx 90.9\%, \qquad
  \text{throughput} = \frac{5\ \text{jobs}}{27.5\ \text{ms}} \approx 0.18\ \text{jobs/ms}.
\]
With zero dispatch latency the same schedule gives $100\%$ and
$5/25 = 0.2$~jobs/ms. Every extra switch moves both numbers down.
\end{example}

\subsection*{Criteria conflict}

Consider a laptop running a text editor, a browser, and a long video export.
The user feels \key{response time}: the editor must react to a keystroke within
milliseconds. A policy that maximized throughput might run the export to
completion first and leave the editor frozen. Criteria conflict, and a
scheduling policy is a choice of which one to favour. ``Maximize throughput''
is the wrong first goal for this machine.

% -----------------------------------------------------------------------------
\section{First-Come, First-Served}
\label{sec:fcfs}

\noindent\emph{Reading:} Silberschatz et al., Ch.~5 \S5.3.1, pp.~188--189
\cite{SilberschatzOSC_operating_system_concepts}; Stallings, Ch.~9 \S9.2,
p.~410 \cite{StallingsOS_internals_design_principles}.

The simplest rule is the one at every ticket counter: serve whoever arrived
first. It is fair-sounding and easy to build --- and it can be very slow.

\begin{definitionbox}{FCFS}
First-Come, First-Served (FCFS) allocates the CPU to the process that requested
it first. The ready queue is a FIFO queue: an arriving PCB joins the tail; the
scheduler takes the head.
\end{definitionbox}

FCFS is \key{non-preemptive}: the running process keeps the CPU until it blocks
or terminates. Its implementation is a linked list with a head and a tail
pointer. Its average waiting time is often long, and it depends heavily on the
order in which processes arrive, as the next two examples show.

\begin{example}[One long job first]
\label{ex:convoy}
\texttt{P1}, \texttt{P2}, \texttt{P3} arrive at time 0 in that order, with
bursts 24, 3, and 3~ms. Under FCFS, \texttt{P1} runs 0--24, \texttt{P2} runs
24--27, and \texttt{P3} runs 27--30 (Figure~\ref{fig:s1}). All arrive at 0, so
each WT is simply the start time: $\mathrm{WT} = 0, 24, 27$. The average WT is
$(0 + 24 + 27)/3 = 17$~ms. Two 3-ms jobs wait behind one 24-ms job.
\end{example}

\begin{figure}[h]
\centering
\begin{tikzpicture}[x=0.4cm,y=1cm]
  % Gantt: S1 under FCFS (generated from the verified simulator); time axis left to right,
  % bars y=0..0.6, time labels below y=0, arrival arrows y=-0.95..-0.5, arrival names below y=-0.95
  \draw[draw=primary-blue, fill=light-bg] (0,0) rectangle (24,0.6);
  \node[font=\scriptsize] at (12,0.3) {\texttt{P1}};
  \draw[draw=primary-blue, fill=light-bg] (24,0) rectangle (27,0.6);
  \node[font=\scriptsize] at (25.5,0.3) {\texttt{P2}};
  \draw[draw=primary-blue, fill=light-bg] (27,0) rectangle (30,0.6);
  \node[font=\scriptsize] at (28.5,0.3) {\texttt{P3}};
  \node[font=\tiny, below] at (0,0) {0};
  \node[font=\tiny, below] at (24,0) {24};
  \node[font=\tiny, below] at (27,0) {27};
  \node[font=\tiny, below] at (30,0) {30};
\end{tikzpicture}
\caption{FCFS with the long job first (Example~\ref{ex:convoy}): average WT
17~ms.}
\label{fig:s1}
\end{figure}

\begin{example}[Same jobs, different arrival order]
\label{ex:reorder}
Now \texttt{P2}, \texttt{P3}, \texttt{P1} arrive at time 0 in that order. FCFS
runs \texttt{P2} 0--3, \texttt{P3} 3--6, \texttt{P1} 6--30
(Figure~\ref{fig:s1b}). The waiting times are $\mathrm{WT}_1 = 6$,
$\mathrm{WT}_2 = 0$, $\mathrm{WT}_3 = 3$, and the average is
$(6+0+3)/3 = 3$~ms --- a factor of $17/3 \approx 5.7$ better than
Example~\ref{ex:convoy}, with the same work and the same finish time of 30~ms.
\end{example}

\begin{figure}[h]
\centering
\begin{tikzpicture}[x=0.4cm,y=1cm]
  % Gantt: S1b under FCFS (generated from the verified simulator); time axis left to right,
  % bars y=0..0.6, time labels below y=0, arrival arrows y=-0.95..-0.5, arrival names below y=-0.95
  \draw[draw=primary-blue, fill=light-bg] (0,0) rectangle (3,0.6);
  \node[font=\scriptsize] at (1.5,0.3) {\texttt{P2}};
  \draw[draw=primary-blue, fill=light-bg] (3,0) rectangle (6,0.6);
  \node[font=\scriptsize] at (4.5,0.3) {\texttt{P3}};
  \draw[draw=primary-blue, fill=light-bg] (6,0) rectangle (30,0.6);
  \node[font=\scriptsize] at (18,0.3) {\texttt{P1}};
  \node[font=\tiny, below] at (0,0) {0};
  \node[font=\tiny, below] at (3,0) {3};
  \node[font=\tiny, below] at (6,0) {6};
  \node[font=\tiny, below] at (30,0) {30};
\end{tikzpicture}
\caption{FCFS with the short jobs first (Example~\ref{ex:reorder}): average WT
3~ms.}
\label{fig:s1b}
\end{figure}

\subsection*{The convoy effect}

Example~\ref{ex:convoy} is the static picture of a dynamic problem. Suppose one
\key{CPU-bound} process holds the CPU for a long burst. Meanwhile many
\key{I/O-bound} processes finish their I/O and queue behind it, and the I/O
devices sit idle. Eventually the CPU-bound process blocks for I/O; the short
jobs run quickly and also block on I/O; now the CPU sits idle. The cycle
repeats, and everyone moves at the speed of the slowest process. This is the
\key{convoy effect}. Its consequence is that FCFS gives lower CPU and device
utilization than letting the short jobs go first; and because it is
non-preemptive, FCFS is unsuitable for time-sharing systems, where every user
needs a share of the CPU at regular intervals
(Silberschatz et al., Ch.~5 \S5.3.1, pp.~188--189).

\begin{example}[The lab server under FCFS]
\label{ex:rfcfs}
Order of arrival is \texttt{P1}, \texttt{P2}, \texttt{P3}, \texttt{P4},
\texttt{P5}, so FCFS runs them in that order back to back: \texttt{P1} 0--6,
\texttt{P2} 6--14, \texttt{P3} 14--16, \texttt{P4} 16--21, \texttt{P5} 21--25
(Figure~\ref{fig:rfcfs}). For example, \texttt{P3} arrives at 2 and finishes at
16, so $\mathrm{TAT}_3 = 14$ and $\mathrm{WT}_3 = 14 - 2 = 12$: it needs only
2~ms but waits 12~ms behind \texttt{P2}'s 8-ms burst. The full table is:

\begin{center}
{\footnotesize\begin{tabular}{lrrrrrr}
\toprule
\textbf{Proc.} & \textbf{AT} & $\boldsymbol{\tau}$ & \textbf{CT} & \textbf{TAT} & \textbf{WT} & \textbf{RT} \\
\midrule
\texttt{P1} & 0 & 6 & 6 & 6 & 0 & 0 \\
\texttt{P2} & 1 & 8 & 14 & 13 & 5 & 5 \\
\texttt{P3} & 2 & 2 & 16 & 14 & 12 & 12 \\
\texttt{P4} & 3 & 5 & 21 & 18 & 13 & 13 \\
\texttt{P5} & 5 & 4 & 25 & 20 & 16 & 16 \\
\midrule
\multicolumn{4}{@{}l}{\textbf{Average}} & 14.2 & 9.2 & 9.2 \\
\bottomrule
\end{tabular}}
\end{center}

\noindent The waiting times sum to $0+5+12+13+16 = 46$, so the average WT is
$46/5 = 9.2$~ms.
\end{example}

\begin{figure}[h]
\centering
\begin{tikzpicture}[x=0.5cm,y=1cm]
  % Gantt: R under FCFS (generated from the verified simulator); time axis left to right,
  % bars y=0..0.6, time labels below y=0, arrival arrows y=-0.95..-0.5, arrival names below y=-0.95
  \draw[draw=primary-blue, fill=light-bg] (0,0) rectangle (6,0.6);
  \node[font=\scriptsize] at (3,0.3) {\texttt{P1}};
  \draw[draw=primary-blue, fill=light-bg] (6,0) rectangle (14,0.6);
  \node[font=\scriptsize] at (10,0.3) {\texttt{P2}};
  \draw[draw=primary-blue, fill=light-bg] (14,0) rectangle (16,0.6);
  \node[font=\scriptsize] at (15,0.3) {\texttt{P3}};
  \draw[draw=primary-blue, fill=light-bg] (16,0) rectangle (21,0.6);
  \node[font=\scriptsize] at (18.5,0.3) {\texttt{P4}};
  \draw[draw=primary-blue, fill=light-bg] (21,0) rectangle (25,0.6);
  \node[font=\scriptsize] at (23,0.3) {\texttt{P5}};
  \node[font=\tiny, below] at (0,0) {0};
  \node[font=\tiny, below] at (6,0) {6};
  \node[font=\tiny, below] at (14,0) {14};
  \node[font=\tiny, below] at (16,0) {16};
  \node[font=\tiny, below] at (21,0) {21};
  \node[font=\tiny, below] at (25,0) {25};
  \draw[-{Latex[length=1.2mm]}, draw=primary-gold] (0,-0.95) -- (0,-0.5);
  \node[font=\tiny, text=primary-gold, below] at (0,-0.95) {P1};
  \draw[-{Latex[length=1.2mm]}, draw=primary-gold] (1,-0.95) -- (1,-0.5);
  \node[font=\tiny, text=primary-gold, below] at (1,-0.95) {P2};
  \draw[-{Latex[length=1.2mm]}, draw=primary-gold] (2,-0.95) -- (2,-0.5);
  \node[font=\tiny, text=primary-gold, below] at (2,-0.95) {P3};
  \draw[-{Latex[length=1.2mm]}, draw=primary-gold] (3,-0.95) -- (3,-0.5);
  \node[font=\tiny, text=primary-gold, below] at (3,-0.95) {P4};
  \draw[-{Latex[length=1.2mm]}, draw=primary-gold] (5,-0.95) -- (5,-0.5);
  \node[font=\tiny, text=primary-gold, below] at (5,-0.95) {P5};
\end{tikzpicture}
\caption{The lab-server jobs under FCFS (Example~\ref{ex:rfcfs}).}
\label{fig:rfcfs}
\end{figure}

% -----------------------------------------------------------------------------
\section{Shortest-Job-First}
\label{sec:sjf}

\noindent\emph{Reading:} Silberschatz et al., Ch.~5 \S5.3.2, pp.~189--192
\cite{SilberschatzOSC_operating_system_concepts}; Stallings, Ch.~9 \S9.2
\cite{StallingsOS_internals_design_principles}.

If order matters so much, pick the order deliberately: run the shortest job
first.

\begin{definitionbox}{SJF}
Shortest-Job-First (SJF) gives the CPU to the ready process with the smallest
next CPU burst. Ties are broken FCFS.
\end{definitionbox}

The more accurate name is \emph{shortest-next-CPU-burst}: the length that
matters is the next burst, not the whole job. In the form we study first, SJF is
\key{non-preemptive}: once chosen, a process runs its whole burst. Its
preemptive form is SRTF (Section~\ref{sec:srtf}). Stallings calls the
non-preemptive policy SPN, \emph{shortest process next}.

\begin{example}[SJF with simultaneous arrivals]
\label{ex:sjf4}
\texttt{P1}--\texttt{P4} arrive at 0 with bursts 6, 8, 7, 3~ms. SJF sorts by
burst: \texttt{P4} (3) runs 0--3, \texttt{P1} (6) runs 3--9, \texttt{P3} (7)
runs 9--16, \texttt{P2} (8) runs 16--24 (Figure~\ref{fig:s2}). The waiting
times are the start times: $\mathrm{WT}_1 = 3$, $\mathrm{WT}_2 = 16$,
$\mathrm{WT}_3 = 9$, $\mathrm{WT}_4 = 0$, average $(3+16+9+0)/4 = 7$~ms. FCFS
in index order would run \texttt{P1}, \texttt{P2}, \texttt{P3}, \texttt{P4}
starting at 0, 6, 14, 21, for an average of $(0+6+14+21)/4 = 10.25$~ms.
\end{example}

\begin{figure}[h]
\centering
\begin{tikzpicture}[x=0.5cm,y=1cm]
  % Gantt: S2 under SJF (generated from the verified simulator); time axis left to right,
  % bars y=0..0.6, time labels below y=0, arrival arrows y=-0.95..-0.5, arrival names below y=-0.95
  \draw[draw=primary-blue, fill=light-bg] (0,0) rectangle (3,0.6);
  \node[font=\scriptsize] at (1.5,0.3) {\texttt{P4}};
  \draw[draw=primary-blue, fill=light-bg] (3,0) rectangle (9,0.6);
  \node[font=\scriptsize] at (6,0.3) {\texttt{P1}};
  \draw[draw=primary-blue, fill=light-bg] (9,0) rectangle (16,0.6);
  \node[font=\scriptsize] at (12.5,0.3) {\texttt{P3}};
  \draw[draw=primary-blue, fill=light-bg] (16,0) rectangle (24,0.6);
  \node[font=\scriptsize] at (20,0.3) {\texttt{P2}};
  \node[font=\tiny, below] at (0,0) {0};
  \node[font=\tiny, below] at (3,0) {3};
  \node[font=\tiny, below] at (9,0) {9};
  \node[font=\tiny, below] at (16,0) {16};
  \node[font=\tiny, below] at (24,0) {24};
\end{tikzpicture}
\caption{SJF on four simultaneous arrivals (Example~\ref{ex:sjf4}): average WT
7~ms, versus 10.25~ms for FCFS.}
\label{fig:s2}
\end{figure}

\subsection*{Why SJF is optimal}

Let $n$ processes all be ready at time 0 and run them in some order with bursts
$\tau_{(1)}, \tau_{(2)}, \ldots, \tau_{(n)}$. The process in position $k$ waits
for the bursts of all processes before it, so the $k$-th burst delays every one
of the $n-k$ processes after it. Summing over processes,
\[
  \sum_{i} \mathrm{WT}_i = \sum_{k=1}^{n} (n-k)\,\tau_{(k)} .
\]
The weights $n-1, n-2, \ldots, 0$ decrease with position, and the largest weight
$n-1$ multiplies the first burst. To minimize the sum, the smallest burst must
take the largest weight, the next smallest the next weight, and so on: that is
exactly SJF order. Equivalently, if two adjacent processes appear long--short,
swapping them into short--long order lowers the total: with the long burst
$L$ in position $k$ and the short burst $S$ in position $k+1$, the swap changes
the sum by $(n-k)S + (n-k-1)L - (n-k)L - (n-k-1)S = S - L < 0$.

\begin{definitionbox}{Optimality of SJF}
For a given set of processes, SJF gives the minimum average waiting time.
\end{definitionbox}

\noindent (Silberschatz et al., Ch.~5 \S5.3.2, from p.~189, states that SJF is
provably optimal.)

\begin{example}[The lab server under non-preemptive SJF]
\label{ex:rsjf}
At $t=0$ only \texttt{P1} is ready, so SJF has no choice: \texttt{P1} runs to
completion at 6. By then \texttt{P2}, \texttt{P3}, \texttt{P4}, \texttt{P5} have
all arrived, with bursts $\{8, 2, 5, 4\}$. SJF runs them shortest first:
\texttt{P3} 6--8, \texttt{P5} 8--12, \texttt{P4} 12--17, \texttt{P2} 17--25
(Figure~\ref{fig:rsjf}).

\begin{center}
{\footnotesize\begin{tabular}{lrrrrrr}
\toprule
\textbf{Proc.} & \textbf{AT} & $\boldsymbol{\tau}$ & \textbf{CT} & \textbf{TAT} & \textbf{WT} & \textbf{RT} \\
\midrule
\texttt{P1} & 0 & 6 & 6 & 6 & 0 & 0 \\
\texttt{P2} & 1 & 8 & 25 & 24 & 16 & 16 \\
\texttt{P3} & 2 & 2 & 8 & 6 & 4 & 4 \\
\texttt{P4} & 3 & 5 & 17 & 14 & 9 & 9 \\
\texttt{P5} & 5 & 4 & 12 & 7 & 3 & 3 \\
\midrule
\multicolumn{4}{@{}l}{\textbf{Average}} & 11.4 & 6.4 & 6.4 \\
\bottomrule
\end{tabular}}
\end{center}

\noindent The waiting times sum to $0+16+4+9+3 = 32$, so the average WT is
$32/5 = 6.4$~ms, against 9.2~ms under FCFS.
\end{example}

\begin{figure}[h]
\centering
\begin{tikzpicture}[x=0.5cm,y=1cm]
  % Gantt: R under SJF (generated from the verified simulator); time axis left to right,
  % bars y=0..0.6, time labels below y=0, arrival arrows y=-0.95..-0.5, arrival names below y=-0.95
  \draw[draw=primary-blue, fill=light-bg] (0,0) rectangle (6,0.6);
  \node[font=\scriptsize] at (3,0.3) {\texttt{P1}};
  \draw[draw=primary-blue, fill=light-bg] (6,0) rectangle (8,0.6);
  \node[font=\scriptsize] at (7,0.3) {\texttt{P3}};
  \draw[draw=primary-blue, fill=light-bg] (8,0) rectangle (12,0.6);
  \node[font=\scriptsize] at (10,0.3) {\texttt{P5}};
  \draw[draw=primary-blue, fill=light-bg] (12,0) rectangle (17,0.6);
  \node[font=\scriptsize] at (14.5,0.3) {\texttt{P4}};
  \draw[draw=primary-blue, fill=light-bg] (17,0) rectangle (25,0.6);
  \node[font=\scriptsize] at (21,0.3) {\texttt{P2}};
  \node[font=\tiny, below] at (0,0) {0};
  \node[font=\tiny, below] at (6,0) {6};
  \node[font=\tiny, below] at (8,0) {8};
  \node[font=\tiny, below] at (12,0) {12};
  \node[font=\tiny, below] at (17,0) {17};
  \node[font=\tiny, below] at (25,0) {25};
  \draw[-{Latex[length=1.2mm]}, draw=primary-gold] (0,-0.95) -- (0,-0.5);
  \node[font=\tiny, text=primary-gold, below] at (0,-0.95) {P1};
  \draw[-{Latex[length=1.2mm]}, draw=primary-gold] (1,-0.95) -- (1,-0.5);
  \node[font=\tiny, text=primary-gold, below] at (1,-0.95) {P2};
  \draw[-{Latex[length=1.2mm]}, draw=primary-gold] (2,-0.95) -- (2,-0.5);
  \node[font=\tiny, text=primary-gold, below] at (2,-0.95) {P3};
  \draw[-{Latex[length=1.2mm]}, draw=primary-gold] (3,-0.95) -- (3,-0.5);
  \node[font=\tiny, text=primary-gold, below] at (3,-0.95) {P4};
  \draw[-{Latex[length=1.2mm]}, draw=primary-gold] (5,-0.95) -- (5,-0.5);
  \node[font=\tiny, text=primary-gold, below] at (5,-0.95) {P5};
\end{tikzpicture}
\caption{The lab-server jobs under non-preemptive SJF (Example~\ref{ex:rsjf}).}
\label{fig:rsjf}
\end{figure}

\subsection*{The catch: predicting the next burst}

The short-term scheduler cannot read the future. It knows how long past bursts
took, but not how long the next one will take, so it can only \emph{predict}
the next CPU burst. The standard predictor is the exponential average.

\begin{definitionbox}{Exponential averaging}
Let $\tau_n$ be the measured length of the $n$-th CPU burst and $\hat{\tau}_n$
its prediction. The prediction for the next burst is
\[
  \hat{\tau}_{n+1} = \alpha\,\tau_n + (1-\alpha)\,\hat{\tau}_n ,
  \qquad 0 \le \alpha \le 1 .
\]
\end{definitionbox}

The parameter $\alpha$ sets how much recent behaviour counts. With $\alpha = 0$,
$\hat{\tau}_{n+1} = \hat{\tau}_n$: only history counts, and recent behaviour is
ignored. With $\alpha = 1$, $\hat{\tau}_{n+1} = \tau_n$: only the last burst
counts. Expanding the recursion,
\[
  \hat{\tau}_{n+1} = \alpha\tau_n + (1-\alpha)\alpha\tau_{n-1} + \cdots
    + (1-\alpha)^j\alpha\tau_{n-j} + \cdots + (1-\alpha)^{n+1}\hat{\tau}_0 ,
\]
so the burst $j$ steps back gets weight $\alpha(1-\alpha)^j$: each step back
counts less. (Silberschatz writes $t_n$ for the measured burst; these notes keep
$\tau$ for bursts and add a hat for a prediction.)

\begin{example}[Predicting bursts with $\alpha = 1/2$]
\label{ex:ema}
Take $\hat{\tau}_0 = 10$~ms and measured bursts $6, 4, 6, 4, 13, 13, 13$~ms. Each
prediction is the average of the last measurement and the last prediction:
\begin{align*}
  \hat{\tau}_1 &= \tfrac12(6) + \tfrac12(10) = 8, &
  \hat{\tau}_2 &= \tfrac12(4) + \tfrac12(8) = 6, \\
  \hat{\tau}_3 &= \tfrac12(6) + \tfrac12(6) = 6, &
  \hat{\tau}_4 &= \tfrac12(4) + \tfrac12(6) = 5, \\
  \hat{\tau}_5 &= \tfrac12(13) + \tfrac12(5) = 9, &
  \hat{\tau}_6 &= \tfrac12(13) + \tfrac12(9) = 11, \\
  \hat{\tau}_7 &= \tfrac12(13) + \tfrac12(11) = 12. &&
\end{align*}

\begin{center}
\begin{tabular}{lcccccccc}
  \toprule
  $n$ & 0 & 1 & 2 & 3 & 4 & 5 & 6 & 7 \\
  \midrule
  measured $\tau_n$ & 6 & 4 & 6 & 4 & 13 & 13 & 13 & --- \\
  predicted $\hat{\tau}_n$ & 10 & 8 & 6 & 6 & 5 & 9 & 11 & 12 \\
  \bottomrule
\end{tabular}
\end{center}

\noindent When the behaviour jumps from 4 to 13, the prediction lags --- 9, 11,
12 --- and converges toward 13. This is the same instance as the
exponential-average illustration in Silberschatz et al., Ch.~5 \S5.3.2.
\end{example}

\begin{example}[Choosing $\alpha$ for an alternating process]
A process alternates bursts of 2~ms and 20~ms. With $\alpha$ near 1 the
predictor follows the last burst: it predicts 2 right before a 20 and 20 right
before a 2 --- always wrong. With $\alpha$ near 0 the prediction settles near
the long-run average, about 11~ms: never exact, but stable. SJF is only as good
as its predictor; a bad guess turns ``shortest first'' into an arbitrary order.
\end{example}

% -----------------------------------------------------------------------------
\section{Shortest-Remaining-Time-First}
\label{sec:srtf}

\noindent\emph{Reading:} Silberschatz et al., Ch.~5 \S5.3.2 (from p.~189)
\cite{SilberschatzOSC_operating_system_concepts}; Stallings, Ch.~9 \S9.2
\cite{StallingsOS_internals_design_principles}.

Non-preemptive SJF ignores a short job that arrives one millisecond after a long
one starts. The preemptive form does not.

\begin{definitionbox}{SRTF}
Shortest-Remaining-Time-First (SRTF) is preemptive SJF. When a process arrives,
its burst is compared with the remaining time of the running process; if it is
strictly shorter, the running process is preempted.
\end{definitionbox}

Decisions now happen at every arrival and every completion (cases 1, 3, and 4
of Section~\ref{sec:concepts}). Equal remaining times do not preempt, and ties
go to the earlier arrival. Waiting time now includes the time a preempted
process spends back in the ready queue, which is why we need the general
formula $\mathrm{WT}_i = \mathrm{TAT}_i - \tau_i$ rather than ``start time minus
arrival''. Stallings calls this policy SRT.

\begin{example}[SRTF versus non-preemptive SJF]
\label{ex:srtf4}
Processes (AT, $\tau$): \texttt{P1} $(0,8)$, \texttt{P2} $(1,4)$, \texttt{P3}
$(2,9)$, \texttt{P4} $(3,5)$. Step by step:
\begin{itemize}
  \item $t=0$: only \texttt{P1}; it starts.
  \item $t=1$: \texttt{P2} arrives with 4; \texttt{P1} has $8-1 = 7$ left.
        $4 < 7$, so \texttt{P1} is preempted and \texttt{P2} runs.
  \item $t=2$: \texttt{P3} arrives with 9; \texttt{P2} has 3 left. No preemption.
  \item $t=3$: \texttt{P4} arrives with 5; \texttt{P2} has 2 left. No preemption.
  \item $t=5$: \texttt{P2} completes. Remaining: \texttt{P1} 7, \texttt{P3} 9,
        \texttt{P4} 5. \texttt{P4} runs 5--10.
  \item $t=10$: \texttt{P1} (7) beats \texttt{P3} (9); \texttt{P1} resumes and
        runs 10--17. Then \texttt{P3} runs 17--26.
\end{itemize}

\begin{center}
{\footnotesize\begin{tabular}{lrrrrrr}
\toprule
\textbf{Proc.} & \textbf{AT} & $\boldsymbol{\tau}$ & \textbf{CT} & \textbf{TAT} & \textbf{WT} & \textbf{RT} \\
\midrule
\texttt{P1} & 0 & 8 & 17 & 17 & 9 & 0 \\
\texttt{P2} & 1 & 4 & 5 & 4 & 0 & 0 \\
\texttt{P3} & 2 & 9 & 26 & 24 & 15 & 15 \\
\texttt{P4} & 3 & 5 & 10 & 7 & 2 & 2 \\
\midrule
\multicolumn{4}{@{}l}{\textbf{Average}} & 13 & 6.5 & 4.25 \\
\bottomrule
\end{tabular}}
\end{center}

\noindent \texttt{P1} waited from 1 to 10, so $\mathrm{WT}_1 = 9$ even though its
$\mathrm{RT}_1 = 0$. The average WT is $(9+0+15+2)/4 = 26/4 = 6.5$~ms. Run
non-preemptively, SJF would execute \texttt{P1} 0--8, \texttt{P2} 8--12,
\texttt{P4} 12--17, \texttt{P3} 17--26, with waits $0, 7, 15, 9$ and average
$31/4 = 7.75$~ms.
\end{example}

\begin{figure}[h]
\centering
\begin{tikzpicture}[x=0.42cm,y=1cm]
  % Gantt: S3 under SRTF (generated from the verified simulator); time axis left to right,
  % bars y=0..0.6, time labels below y=0, arrival arrows y=-0.95..-0.5, arrival names below y=-0.95
  \draw[draw=primary-blue, fill=light-bg] (0,0) rectangle (1,0.6);
  \node[font=\scriptsize] at (0.5,0.3) {\texttt{P1}};
  \draw[draw=primary-blue, fill=light-bg] (1,0) rectangle (5,0.6);
  \node[font=\scriptsize] at (3,0.3) {\texttt{P2}};
  \draw[draw=primary-blue, fill=light-bg] (5,0) rectangle (10,0.6);
  \node[font=\scriptsize] at (7.5,0.3) {\texttt{P4}};
  \draw[draw=primary-blue, fill=light-bg] (10,0) rectangle (17,0.6);
  \node[font=\scriptsize] at (13.5,0.3) {\texttt{P1}};
  \draw[draw=primary-blue, fill=light-bg] (17,0) rectangle (26,0.6);
  \node[font=\scriptsize] at (21.5,0.3) {\texttt{P3}};
  \node[font=\tiny, below] at (0,0) {0};
  \node[font=\tiny, below] at (1,0) {1};
  \node[font=\tiny, below] at (5,0) {5};
  \node[font=\tiny, below] at (10,0) {10};
  \node[font=\tiny, below] at (17,0) {17};
  \node[font=\tiny, below] at (26,0) {26};
  \draw[-{Latex[length=1.2mm]}, draw=primary-gold] (0,-0.95) -- (0,-0.5);
  \node[font=\tiny, text=primary-gold, below] at (0,-0.95) {P1};
  \draw[-{Latex[length=1.2mm]}, draw=primary-gold] (1,-0.95) -- (1,-0.5);
  \node[font=\tiny, text=primary-gold, below] at (1,-0.95) {P2};
  \draw[-{Latex[length=1.2mm]}, draw=primary-gold] (2,-0.95) -- (2,-0.5);
  \node[font=\tiny, text=primary-gold, below] at (2,-0.95) {P3};
  \draw[-{Latex[length=1.2mm]}, draw=primary-gold] (3,-0.95) -- (3,-0.5);
  \node[font=\tiny, text=primary-gold, below] at (3,-0.95) {P4};
\end{tikzpicture}
\caption{SRTF schedule of Example~\ref{ex:srtf4}: \texttt{P1} is preempted at 1
and resumes at 10.}
\label{fig:s3}
\end{figure}

\begin{example}[The lab server under SRTF]
\label{ex:rsrtf}
\texttt{P1} starts at 0. At $t=1$, \texttt{P2} (8) does not beat \texttt{P1}'s
remaining 5. At $t=2$, \texttt{P3} (2) beats \texttt{P1}'s remaining 4:
\texttt{P1} is preempted and \texttt{P3} runs 2--4. At $t=3$, \texttt{P4} (5)
does not beat \texttt{P3}'s remaining 1. At $t=4$, \texttt{P3} completes; the
remaining times are \texttt{P1} 4, \texttt{P2} 8, \texttt{P4} 5, so \texttt{P1}
resumes. At $t=5$, \texttt{P5} (4) does not beat \texttt{P1}'s remaining 3, so
\texttt{P1} finishes at 8. Then \texttt{P5} (4) runs 8--12, \texttt{P4} (5)
12--17, and \texttt{P2} (8) 17--25 (Figure~\ref{fig:rsrtf}).

\begin{center}
{\footnotesize\begin{tabular}{lrrrrrr}
\toprule
\textbf{Proc.} & \textbf{AT} & $\boldsymbol{\tau}$ & \textbf{CT} & \textbf{TAT} & \textbf{WT} & \textbf{RT} \\
\midrule
\texttt{P1} & 0 & 6 & 8 & 8 & 2 & 0 \\
\texttt{P2} & 1 & 8 & 25 & 24 & 16 & 16 \\
\texttt{P3} & 2 & 2 & 4 & 2 & 0 & 0 \\
\texttt{P4} & 3 & 5 & 17 & 14 & 9 & 9 \\
\texttt{P5} & 5 & 4 & 12 & 7 & 3 & 3 \\
\midrule
\multicolumn{4}{@{}l}{\textbf{Average}} & 11 & 6 & 5.6 \\
\bottomrule
\end{tabular}}
\end{center}

\noindent The waiting times sum to $2+16+0+9+3 = 30$: average WT
$30/5 = 6.0$~ms. The response times sum to 28: average RT 5.6~ms.
\end{example}

\begin{figure}[h]
\centering
\begin{tikzpicture}[x=0.5cm,y=1cm]
  % Gantt: R under SRTF (generated from the verified simulator); time axis left to right,
  % bars y=0..0.6, time labels below y=0, arrival arrows y=-0.95..-0.5, arrival names below y=-0.95
  \draw[draw=primary-blue, fill=light-bg] (0,0) rectangle (2,0.6);
  \node[font=\scriptsize] at (1,0.3) {\texttt{P1}};
  \draw[draw=primary-blue, fill=light-bg] (2,0) rectangle (4,0.6);
  \node[font=\scriptsize] at (3,0.3) {\texttt{P3}};
  \draw[draw=primary-blue, fill=light-bg] (4,0) rectangle (8,0.6);
  \node[font=\scriptsize] at (6,0.3) {\texttt{P1}};
  \draw[draw=primary-blue, fill=light-bg] (8,0) rectangle (12,0.6);
  \node[font=\scriptsize] at (10,0.3) {\texttt{P5}};
  \draw[draw=primary-blue, fill=light-bg] (12,0) rectangle (17,0.6);
  \node[font=\scriptsize] at (14.5,0.3) {\texttt{P4}};
  \draw[draw=primary-blue, fill=light-bg] (17,0) rectangle (25,0.6);
  \node[font=\scriptsize] at (21,0.3) {\texttt{P2}};
  \node[font=\tiny, below] at (0,0) {0};
  \node[font=\tiny, below] at (2,0) {2};
  \node[font=\tiny, below] at (4,0) {4};
  \node[font=\tiny, below] at (8,0) {8};
  \node[font=\tiny, below] at (12,0) {12};
  \node[font=\tiny, below] at (17,0) {17};
  \node[font=\tiny, below] at (25,0) {25};
  \draw[-{Latex[length=1.2mm]}, draw=primary-gold] (0,-0.95) -- (0,-0.5);
  \node[font=\tiny, text=primary-gold, below] at (0,-0.95) {P1};
  \draw[-{Latex[length=1.2mm]}, draw=primary-gold] (1,-0.95) -- (1,-0.5);
  \node[font=\tiny, text=primary-gold, below] at (1,-0.95) {P2};
  \draw[-{Latex[length=1.2mm]}, draw=primary-gold] (2,-0.95) -- (2,-0.5);
  \node[font=\tiny, text=primary-gold, below] at (2,-0.95) {P3};
  \draw[-{Latex[length=1.2mm]}, draw=primary-gold] (3,-0.95) -- (3,-0.5);
  \node[font=\tiny, text=primary-gold, below] at (3,-0.95) {P4};
  \draw[-{Latex[length=1.2mm]}, draw=primary-gold] (5,-0.95) -- (5,-0.5);
  \node[font=\tiny, text=primary-gold, below] at (5,-0.95) {P5};
\end{tikzpicture}
\caption{The lab-server jobs under SRTF (Example~\ref{ex:rsrtf}).}
\label{fig:rsrtf}
\end{figure}

% -----------------------------------------------------------------------------
\section{Comparing the Three Policies}
\label{sec:compare}

Table~\ref{tab:compare} collects the averages of
Examples~\ref{ex:rfcfs}, \ref{ex:rsjf}, and \ref{ex:rsrtf}.

\begin{table}[h]
\centering
\renewcommand{\arraystretch}{1.2}
\begin{tabular}{lcccc}
  \toprule
  \textbf{Policy} & \textbf{avg WT} & \textbf{avg TAT} & \textbf{avg RT} & \textbf{dispatches} \\
  \midrule
  FCFS & 9.2 & 14.2 & 9.2 & 5 \\
  SJF (non-preemptive) & 6.4 & 11.4 & 6.4 & 5 \\
  SRTF & 6.0 & 11.0 & 5.6 & 6 \\
  \bottomrule
\end{tabular}
\caption{Three policies on the lab-server workload (all times in ms).}
\label{tab:compare}
\end{table}

All three schedules finish at 25~ms: the work is the same, and only the
\emph{order} --- and hence the waiting --- changes. SRTF buys the best averages
with one extra dispatch (its Gantt chart has six bars). With the 0.5-ms dispatch
latency of Example~\ref{ex:util}, SRTF's utilization is
$25/(25 + 6 \times 0.5) = 25/28 \approx 89.3\%$, against $90.9\%$ for FCFS.

\subsection*{The price of short-first: starvation}

If short jobs keep arriving, a long job can wait indefinitely: SJF and SRTF can
\key{starve} long processes. On the lab server, \texttt{P2} (8~ms) already has
the largest WT (16~ms) under both SJF and SRTF; a steady stream of 2-ms jobs
would postpone it forever. FCFS never starves anyone, because each process waits
only for those ahead of it in the queue.

SJF and SRTF carry two further costs: they need a burst prediction, and SRTF
pays a dispatch on every preemption. Ageing --- raising a process's priority
with the time it has waited --- is the usual cure for starvation, and is taken
up in Lecture~9.

\subsection*{Common misconceptions}

\begin{itemize}
  \item \bad{``Waiting time is completion minus arrival.''} That is TAT;
        $\mathrm{WT} = \mathrm{TAT} - \tau$.
  \item \bad{``SJF uses the total length of the job.''} It uses the predicted
        \emph{next CPU burst}.
  \item \bad{``SRTF preempts on equal remaining time.''} Only a strictly shorter
        newcomer preempts.
  \item \bad{``A better policy finishes all the work sooner.''} The last
        completion is the same; the averages change.
\end{itemize}

% -----------------------------------------------------------------------------
\section{Summary}

\begin{itemize}
  \item Processes alternate \key{CPU bursts} and \key{I/O bursts}; the
        short-term scheduler fills the CPU while others wait for I/O.
  \item Scheduling happens on block and exit (non-preemptive) and also on timer
        or I/O completion (\key{preemptive}); the \key{dispatcher} pays a latency
        on every switch.
  \item $\mathrm{TAT} = \mathrm{CT} - \mathrm{AT}$,
        $\mathrm{WT} = \mathrm{TAT} - \tau$,
        $\mathrm{RT} = \mathrm{ST} - \mathrm{AT}$.
  \item \key{FCFS} is simple and starvation-free but suffers the convoy effect.
  \item \key{SJF} minimizes average WT but needs a predicted burst (exponential
        average); \key{SRTF} is its preemptive form. Both can starve long jobs.
\end{itemize}

\section{Looking Ahead and Further Reading}

This chapter covered what scheduling is, how to measure it, and the two policies
that order by arrival and by burst length. The next lecture covers
\key{Priority} scheduling and ageing, \key{Round Robin} with a time quantum,
multilevel queues and feedback queues, and scheduling on more than one CPU.
Reading for it: Silberschatz et al., Ch.~5 \S5.3--5.5
\cite{SilberschatzOSC_operating_system_concepts}; Stallings, Ch.~9 \S9.2 and
Ch.~10 \S10.1 \cite{StallingsOS_internals_design_principles}.

\begin{itemize}
  \item A.~Silberschatz, P.~B.~Galvin, G.~Gagne, \emph{Operating System
        Concepts}, 8e --- Ch.~5 \S5.1 (basic concepts, p.~183), \S5.2
        (scheduling criteria, p.~187), \S5.3.1 (FCFS, p.~188), \S5.3.2 (SJF and
        SRTF, p.~189).
  \item W.~Stallings, \emph{Operating Systems: Internals and Design
        Principles}, 6e --- Ch.~9 \S9.1 (types of scheduling, p.~406), \S9.2
        (scheduling algorithms: FCFS, SPN, SRT, p.~410).
\end{itemize}

% -----------------------------------------------------------------------------
\section*{Self-Check}
\addcontentsline{toc}{section}{Self-Check}

\subsection*{Exercises}
\begin{enumerate}
  \item Processes (AT, $\tau$) in ms: \texttt{P1} $(0,3)$, \texttt{P2} $(1,5)$,
        \texttt{P3} $(2,1)$, \texttt{P4} $(3,2)$. Draw the FCFS and the
        non-preemptive SJF Gantt charts and give the average WT of each.
  \item Processes (AT, $\tau$): \texttt{P1} $(0,5)$, \texttt{P2} $(1,1)$,
        \texttt{P3} $(2,3)$. Compute the average WT under non-preemptive SJF and
        under SRTF, and the number of dispatches SRTF makes.
  \item With $\alpha = 0.25$ and $\hat{\tau}_0 = 8$~ms, the first three
        measured bursts are 4, 12, 6~ms. Compute $\hat{\tau}_1$,
        $\hat{\tau}_2$, $\hat{\tau}_3$.
  \item Three processes arrive at 0 with bursts 5, 2, 1~ms. Using the formula
        $\sum_i \mathrm{WT}_i = \sum_k (n-k)\tau_{(k)}$, compute the average WT
        for the order 5, 2, 1 and for SJF order, and explain why no order can do
        better than SJF.
  \item Explain why, for a non-preemptive policy, $\mathrm{WT}_i =
        \mathrm{RT}_i$, and give a one-line counterexample under SRTF from
        Example~\ref{ex:srtf4}.
\end{enumerate}

\subsection*{Solutions}
\begingroup\small
\begin{enumerate}
  \item \emph{FCFS:} \texttt{P1} 0--3, \texttt{P2} 3--8, \texttt{P3} 8--9,
        \texttt{P4} 9--11. WT $= 0, 2, 6, 6$; average $14/4 = 3.5$~ms.
        \emph{SJF:} at 0 only \texttt{P1}, which runs 0--3. At 3 the ready
        bursts are \texttt{P2} 5, \texttt{P3} 1, \texttt{P4} 2: \texttt{P3}
        3--4, \texttt{P4} 4--6, \texttt{P2} 6--11. WT: \texttt{P1} 0,
        \texttt{P2} $6-1 = 5$, \texttt{P3} $3-2 = 1$, \texttt{P4} $4-3 = 1$;
        average $7/4 = 1.75$~ms.
  \item \emph{SJF:} \texttt{P1} runs 0--5 (only process at 0); at 5,
        \texttt{P2} (1) then \texttt{P3} (3): \texttt{P2} 5--6, \texttt{P3}
        6--9. WT $= 0, 4, 4$; average $8/3 \approx 2.67$~ms.
        \emph{SRTF:} at 1, \texttt{P2} (1) beats \texttt{P1}'s remaining 4 and
        runs 1--2; at 2, \texttt{P3} (3) beats \texttt{P1}'s remaining 4 and runs
        2--5; \texttt{P1} resumes 5--9. CT $= 9, 2, 5$; TAT $= 9, 1, 3$;
        WT $= 4, 0, 0$; average $4/3 \approx 1.33$~ms. SRTF makes 4 dispatches
        (\texttt{P1}, \texttt{P2}, \texttt{P3}, \texttt{P1}).
  \item $\hat{\tau}_1 = 0.25(4) + 0.75(8) = 7$;
        $\hat{\tau}_2 = 0.25(12) + 0.75(7) = 8.25$;
        $\hat{\tau}_3 = 0.25(6) + 0.75(8.25) = 7.6875$~ms.
  \item Order 5, 2, 1: $\sum \mathrm{WT} = 2(5) + 1(2) + 0(1) = 12$, average
        $4$~ms. SJF order 1, 2, 5: $2(1) + 1(2) + 0(5) = 4$, average
        $4/3 \approx 1.33$~ms. The weights $2, 1, 0$ decrease with position, so
        the sum is smallest when the smallest burst takes the largest weight;
        any long--short adjacent pair can be swapped to lower the sum by
        $L - S > 0$, so only the sorted order cannot be improved.
  \item Non-preemptively, process $i$ runs once from $\mathrm{ST}_i$ to
        $\mathrm{CT}_i = \mathrm{ST}_i + \tau_i$, so
        $\mathrm{WT}_i = \mathrm{CT}_i - \mathrm{AT}_i - \tau_i =
        \mathrm{ST}_i - \mathrm{AT}_i = \mathrm{RT}_i$. Under SRTF in
        Example~\ref{ex:srtf4}, \texttt{P1} has $\mathrm{RT}_1 = 0$ but
        $\mathrm{WT}_1 = 9$, because it waited from 1 to 10 after preemption.
\end{enumerate}
\endgroup

\bibliographystyle{plain}
\bibliography{../../Bibliography_base}

\end{document}
